\documentclass[11pt,a4paper]{article}
\usepackage[margin=1in]{geometry}
\usepackage[utf8]{inputenc}
\usepackage[T1]{fontenc}
\usepackage{amsmath,amssymb,amsthm}
\usepackage{graphicx}
\usepackage{booktabs}
\usepackage{longtable}
\usepackage{parskip}
\usepackage{xcolor}
\usepackage[numbers,sort&compress]{natbib}
\usepackage{hyperref}
\hypersetup{colorlinks=true,linkcolor=blue,urlcolor=cyan,citecolor=red}

\newtheorem{theorem}{Theorem}[section]
\newtheorem{proposition}[theorem]{Proposition}
\newtheorem{corollary}[theorem]{Corollary}
\newtheorem{remark}[theorem]{Remark}
\newtheorem{definition}[theorem]{Definition}

\newcommand{\lP}{l_P}
\newcommand{\tP}{t_P}
\newcommand{\EP}{E_P}
\newcommand{\mP}{m_P}
\newcommand{\SSV}{\mathrm{SSV}}
\newcommand{\ket}[1]{|#1\rangle}
\newcommand{\bra}[1]{\langle#1|}

\title{\textbf{Quantum Mechanics from a Planck-Scale Lattice:\\
Five Theorems in Conscious Point Physics}\\[8pt]
{\normalsize Prepared for submission to \textit{Foundations of Physics}}}

\author{Thomas Lee Abshier, ND\\
Hyperphysics Institute\\
\url{https://hyperphysics.com}\\
\texttt{drthomas007@protonmail.com}}

\date{22 March 2026}

\begin{document}
\maketitle

\begin{abstract}
We present five theorems that together constitute a derivation of
non-relativistic quantum mechanics and the free quantum field from
a discrete Planck-scale lattice theory — Conscious Point Physics (CPP).
The primitives are: Conscious Points (CPs) carrying charge $\pm e$;
Displacement Increment (DI) bits carrying both number density $\rho$ and
geometric phase $\phi$, with complex amplitude
$\psi_i = \sqrt{\rho_i}\,e^{i\phi_i}$ at each Grid Point (GP) of the
lattice; a 600-cell lattice with 120 GPs, coordination $z=12$, and
H$_4$ icosahedral symmetry; and the Nexus, an atemporal global
conservation law enforcing DI-bit number and phase circulation at each
Absolute Moment tick.

\textbf{Theorem 1} (Schr\"{o}dinger): Complex hopping on the 600-cell
with amplitude $T = \hbar^2/(4m\Delta s^2)$ gives
$i\hbar\partial_t\psi = [-\hbar^2\nabla^2/(2m) + V]\psi$ in the
continuum limit $\Delta s \to 0$.

\textbf{Theorem 2} (Born rule): The detection probability equals the
DI-bit number density, $P = \rho_{\rm bit} = |\psi|^2$; the
derivation is non-circular because $\psi$ and $A_0$ are defined
independently of probability.

\textbf{Theorem 3} (Bell -- CHSH $= 2\sqrt{2}$): The singlet state is
non-separable; the Nexus maintains it globally; the Born rule gives
$E(\hat{a},\hat{b}) = -\cos\theta$, yielding $|S| = 2\sqrt{2}$ and
exact no-signaling.

\textbf{Theorem 4} (Lindblad -- pointer basis): DP~Sea scattering
under the Born--Markov approximation yields the Lindblad master
equation; the pointer basis is proved geometrically to be the SSV
eigenstates, fixed by the lattice rather than by the apparatus.

\textbf{Theorem 5} (Commutation -- fermion/boson): Bosonic
commutation relations $[a_k, a_{k'}^\dagger] = \delta_{kk'}$ follow
from eigenmode orthonormality of the 600-cell adjacency matrix;
fermion--boson statistics follow from CP exclusion vs.\ neutral
DI-bit modes.

All lattice corrections are of order $(\lP/\lambda)^2 \lesssim 10^{-66}$
at laboratory scales and unobservable. The hierarchy problem becomes
finite fine-tuning rather than infinite renormalization. SM particle
masses and $\alpha = 1/137$ remain open problems.
\end{abstract}

\tableofcontents
\newpage

%=====================================================================
\section{Introduction and Motivation}
%=====================================================================

The five fundamental postulates of standard quantum mechanics — the
Hilbert space, the Born rule, linear superposition, wavefunction collapse,
and field commutation relations — are independent axioms with no
mechanical derivation. Conscious Point Physics (CPP) attempts to derive
each from the dynamics of a single discrete structure: phase-carrying
DI bits propagating on a Planck-scale 600-cell lattice under global
conservation by the Nexus.

This paper is a compressed synthesis of the CPP quantum mechanics series
(Papers 2--6, \citep{cpp2040a,cpp2040b,cpp2040c,cpp2040d,cpp2040e}).
It states and proves the five key theorems, identifies the specific
physical ingredient at each step, and clearly flags what remains open.
Readers seeking extended derivations, comparison with standard theory,
and numerical predictions are referred to the individual series papers.

\subsection{The four primitives}

\begin{definition}[CPP primitives]
\label{def:primitives}
\hfill
\begin{enumerate}
\item \textbf{Conscious Points (CPs):} Planck-scale entities each
  carrying charge $\pm e$ (or $\pm e/3$, $\pm 2e/3$ for quark
  aggregates \citep{cpp5014}). Each occupies one Grid Point.
\item \textbf{DI bits:} Relational quanta exchanged between CPs. Each
  DI bit carries density $\rho \ge 0$ and geometric phase $\phi$,
  forming the complex amplitude $\psi = \sqrt{\rho}\,e^{i\phi}$. DI
  bits propagate at $c = \lP/\tP$.
\item \textbf{600-cell lattice:} 120 Grid Points, 720 edges,
  coordination $z = 12$. Adjacency matrix $A$ (real symmetric,
  $120\times120$) has exactly six distinct eigenvalues
  $\lambda \in \{12,\, 1+\varphi,\, \varphi-1,\, 1-\varphi,\,
  -\varphi,\, -(1+\varphi)\}$ where $\varphi = (1+\sqrt5)/2$
  \citep{brouwer1989}.
\item \textbf{Nexus:} Atemporal global constraint that enforces
  (i)~DI-bit number conservation and (ii)~total phase circulation
  across all 120 Grid Points at every Absolute Moment tick $\Delta t
  = \tP$. It is not a signal; it is not a local hidden variable.
\end{enumerate}
\end{definition}

%=====================================================================
\section{Theorem 1: Schr\"{o}dinger Equation from 600-Cell Hopping}
%=====================================================================

\begin{theorem}[Schr\"{o}dinger equation]
\label{thm:schrodinger}
Let the DI-bit complex amplitude $\psi_i(t)$ at Grid Point $i$ evolve
by complex hopping:
\begin{equation}
\psi_i(t+\Delta t) = \psi_i(t)
- \frac{i\Delta t}{\hbar}\sum_j H_{ij}\psi_j(t),
\label{eq:lattice_evol}
\end{equation}
where
\begin{equation}
H_{ij} = \begin{cases}
-T & j \sim i \ \text{(nearest neighbour)}, \\
+zT & j = i, \quad z = 12, \\
0 & \text{otherwise},
\end{cases}
\qquad T = \frac{\hbar^2}{4m\Delta s^2}.
\label{eq:H}
\end{equation}
In the continuum limit $\Delta s \to 0$ (with $T\Delta s^2 =
\hbar^2/(4m)$ fixed), equation~\eqref{eq:lattice_evol} becomes
\begin{equation}
i\hbar\,\frac{\partial\psi}{\partial t}
= -\frac{\hbar^2}{2m}\nabla^2\psi + V(\mathbf{r})\psi,
\label{eq:Schrodinger}
\end{equation}
where $V(\mathbf{r}) = -k_{\rm PSR}\Delta|\SSV(\mathbf{r})|$ is the
external potential from the Space Stress Vector field.
\end{theorem}

\begin{proof}
The 600-cell graph Laplacian satisfies (Appendix~\ref{app:laplacian}):
\begin{equation}
(\mathcal{A}\psi)_i \equiv \sum_{j\sim i}(\psi_j - \psi_i)
= 2\Delta s^2\nabla^2\psi + O(\Delta s^4).
\label{eq:laplacian}
\end{equation}
Substituting into~\eqref{eq:lattice_evol}:
$i\hbar(\psi_i(t+\Delta t) - \psi_i)/\Delta t = -T\cdot(\mathcal{A}\psi)_i
+ V_i\psi_i = -[\hbar^2/(4m\Delta s^2)]\cdot 2\Delta s^2\nabla^2\psi
+ V\psi$.
Taking $\Delta t \to 0$ yields~\eqref{eq:Schrodinger}. \qed
\end{proof}

\begin{remark}[Why $T = \hbar^2/(4m\Delta s^2)$, not $\hbar^2/(2m\Delta s^2)$]
The 600-cell has $z=12$ in $d=3$ spatial dimensions.
The graph Laplacian sum yields a prefactor $z/(2d) = 12/6 = 2$,
so $(\mathcal{A}\psi)/\Delta s^2 \to 2\nabla^2\psi$.
Setting $T\cdot 2 = \hbar^2/(2m)$ gives $T = \hbar^2/(4m\Delta s^2)$.
Using $T = \hbar^2/(2m\Delta s^2)$ would give $-\hbar^2\nabla^2\psi/m$
(wrong by a factor of~2).
\end{remark}

\begin{remark}[Why the imaginary unit arises from DI-bit phase]
Classical diffusion replaces $-i\Delta t/\hbar$ in~\eqref{eq:lattice_evol}
with $+D\Delta t$ (real, positive), giving the Fokker--Planck equation
$\partial_t\rho = D\nabla^2\rho$ — real diffusion, no wave behaviour.
The key CPP physical ingredient is that each DI-bit hop accumulates phase:
$\Delta\phi = m_{\rm CP}c\Delta s/\hbar$. This geometric phase
accumulation inserts the factor $-i/\hbar$ into the hopping, converting
real diffusion into complex quantum evolution. The Fokker--Planck route
was explicitly ruled out in CPP-2040a~\citep{cpp2040a}: no continuum
limit of a real diffusion equation can produce the imaginary unit.
\end{remark}

%=====================================================================
\section{Theorem 2: Born Rule from DI-Bit Density}
%=====================================================================

\begin{theorem}[Born rule]
\label{thm:born}
The probability of detecting a CP state change at Grid Point $i$ is
$P(i) = |\psi(i)|^2$.
\end{theorem}

\begin{proof}
From companion~C3~\citep{abshier2026c3}: the probability of a CP state
change at a Grid Point is proportional to the local DI-bit number
density, $P \propto \rho_{\rm bit}$.
Since $\psi = \sqrt{\rho_{\rm bit}}\,e^{i\phi}$, we have
$\rho_{\rm bit} = |\psi|^2$ and therefore $P(i) \propto |\psi(i)|^2$. \qed
\end{proof}

\begin{remark}[Non-circularity]
The complex amplitude $\psi$ and the emission amplitude $A_0$ entering
the path sum $\psi(d) = \sum_k A_0 e^{i\phi_k}$ are defined geometrically
from lattice structure and DI-bit density — not from any probability.
The Born rule is a consequence, not a postulate.
\end{remark}

%=====================================================================
\section{Theorem 3: Bell--CHSH $= 2\sqrt{2}$ from the Nexus}
%=====================================================================

\begin{theorem}[Non-separability of the singlet]
\label{thm:nonsep}
The singlet state $\ket{\Psi^-} = (\ket{\uparrow\downarrow} -
\ket{\downarrow\uparrow})/\sqrt{2}$ cannot be written as
$\ket{\phi_A}\otimes\ket{\phi_B}$ for any single-particle states.
\end{theorem}

\begin{proof}
Assume the factorisation. Expanding and matching the four coefficients
$\alpha_A\gamma_B = 0$, $\alpha_A\delta_B = 1/\sqrt2$,
$\beta_A\gamma_B = -1/\sqrt2$, $\beta_A\delta_B = 0$ yields a
contradiction: the first and fourth equations force $\alpha_A\delta_B = 0$
or $\beta_A\gamma_B = 0$, contradicting the middle two. \qed
\end{proof}

\begin{remark}[Role of the Nexus]
The Nexus maintains the joint non-separable state~\eqref{eq:singlet}
globally at every tick, regardless of the spatial separation of particles
A and B. This is not a signal (no information is transmitted) and not a
local hidden variable (it does not factorise, so Bell's theorem does not
apply to it).
\end{remark}

\begin{theorem}[CHSH $= 2\sqrt{2}$]
\label{thm:chsh}
For the singlet state with the Born rule (Theorem~\ref{thm:born}),
the singlet correlation is $E(\hat{a},\hat{b}) = -\cos\theta$,
yielding $|S_{\rm CHSH}| = 2\sqrt{2}$ at optimal settings
($\hat{a}=0°$, $\hat{a}'=90°$, $\hat{b}=45°$, $\hat{b}'=-45°$).
\end{theorem}

\begin{proof}
The joint probabilities from the Born rule for the singlet are
$P(\pm,\pm) = (1\mp\cos\theta)/4$.
The correlation is
$E = P(++) + P(--) - P(+-) - P(-+)
= [(1{-}c) + (1{-}c) - (1{+}c) - (1{+}c)]/4 = -\cos\theta$.
At optimal angles all four $|\theta| = 45°$ except $\theta_{a'b'} = 135°$:
$S = -1/\sqrt2 - 1/\sqrt2 - 1/\sqrt2 - (1/\sqrt2) = -2\sqrt2$.
$|S| = 2\sqrt2 > 2$ violates the classical Bell bound. \qed
\end{proof}

\begin{corollary}[No-signaling]
Alice's marginal probability $P(A{=}+1) = \tfrac12$ is independent of
Bob's measurement axis; symmetrically for Bob.
\end{corollary}

%=====================================================================
\section{Theorem 4: Lindblad Equation and Lattice-Fixed Pointer Basis}
%=====================================================================

\begin{theorem}[Lindblad from DP~Sea scattering]
\label{thm:lindblad}
The reduced density matrix of a DI-bit qubit coupled to the DP~Sea via
$H_{\rm int} = \sum_j g_j(a_j + a_j^\dagger)\otimes\hat{\sigma}_z$
evolves, under the Born--Markov approximation, as
\begin{equation}
\frac{d\rho}{dt} = -\frac{i}{\hbar}[H_S,\rho]
+ \gamma(\hat{\sigma}_z\rho\hat{\sigma}_z - \rho),
\qquad
\gamma = \frac{(\text{sea\_strength})^2\,\EP}{\hbar},
\label{eq:lindblad}
\end{equation}
with sea\_strength $= 0.185$ (companion~C1~\citep{abshier2026c1}).
Off-diagonal elements decay as
$\rho_{01}(t) = \rho_{01}(0)e^{-2\gamma t}$; populations are
conserved.
\end{theorem}

\begin{proof}[Derivation sketch]
The Born--Markov second-order time-convolutionless expansion of the
interaction-picture master equation gives
$\dot\rho_S = -(1/\hbar^2)\int_0^\infty d\tau\,
\mathrm{tr}_B[H_{\rm int}(t),[H_{\rm int}(t-\tau),\rho_S\otimes\rho_B]]$.
With bath correlation time $\tau_{\rm corr} = \tP$ (shortest time in the
theory) and weak coupling sea\_strength $\ll 1$, the Markov approximation
collapses the bath correlator to a delta function, yielding the
Lindblad form~\eqref{eq:lindblad} with $\gamma$ set by the spectral
density of the DP~Sea at zero frequency. \qed
\end{proof}

\begin{theorem}[Pointer basis $=$ SSV eigenstates]
\label{thm:pointer}
The states robust under DP~Sea coupling are the eigenstates of
$\hat{\sigma}_z$, i.e., the DI-bit states of definite phase
projection onto the local SSV direction.
\end{theorem}

\begin{proof}
A state is robust under repeated dephasing kicks proportional to
$\hat{\sigma}_z$ if and only if it commutes with $\hat{\sigma}_z$.
The eigenstates of $\hat{\sigma}_z$ satisfy this; all superpositions
accumulate random phase offsets that drive $\rho_{01} \to 0$.
The 12-edge broadcast at each Grid Point selects the dominant SSV
component, so $\hat{\sigma}_z$ is the SSV projection operator, and
the pointer basis is therefore the SSV eigenstates. \qed
\end{proof}

\begin{remark}[Stronger than einselection]
Standard einselection \citep{zurek1981} identifies the pointer basis
with the eigenstates of the system--environment coupling, which is
apparatus-dependent. In CPP, the coupling is always $\hat{\sigma}_z$
(the SSV projection), so the pointer basis is fixed by the lattice
geometry, independent of the measurement apparatus.
This is a falsifiable CPP-specific prediction.
\end{remark}

%=====================================================================
\section{Theorem 5: Commutation Relations and Fermion--Boson Statistics}
%=====================================================================

\begin{theorem}[Bosonic commutation relations]
\label{thm:commutators}
Define the eigenmode annihilation operator
$a_k = \sum_i u_k(i)^*\hat{c}_i$
where $\mathbf{u}_k$ are the orthonormal eigenvectors of the 600-cell
adjacency matrix $A$ and $\hat{c}_i$ annihilates a DI bit at site $i$.
If site operators satisfy $[\hat{c}_i,\hat{c}_j^\dagger] = \delta_{ij}$,
then
\begin{equation}
[a_k, a_{k'}^\dagger] = \delta_{kk'}.
\label{eq:commutators}
\end{equation}
\end{theorem}

\begin{proof}
$[a_k,a_{k'}^\dagger] = \sum_{i,j}u_k(i)^*u_{k'}(j)[\hat{c}_i,\hat{c}_j^\dagger]
= \sum_i u_k(i)^*u_{k'}(i) = \mathbf{u}_k^\dagger\mathbf{u}_{k'}
= \delta_{kk'}$. \qed
\end{proof}

\begin{theorem}[Fermion--boson distinction]
\label{thm:fermboson}
Charged CP aggregates obey the CP exclusion principle (at most one per
Grid Point), yielding fermionic anticommutators. Neutral DI-bit modes
have no occupancy restriction, yielding bosonic commutators.
\end{theorem}

\begin{proof}
Two same-sign charged CPs at one Grid Point would have zero separation
and therefore infinite SSV self-repulsion (CPP-5014~\citep{cpp5014}).
The site algebra for charged modes therefore has $\hat{c}_i^2 = 0$
(Pauli), giving anticommutators.
Neutral DI bits carry no charge and experience no such exclusion;
their site algebra gives commutators. \qed
\end{proof}

\begin{remark}[Finite renormalization]
The physical Planck-scale lattice provides a UV cutoff
$k_{\max} = \pi/\lP$, rendering every Feynman loop integral finite.
The electron self-energy correction is $\Sigma \sim e^2\EP^2/(4\pi^2)$
— large but finite. The hierarchy problem becomes a well-defined
finite fine-tuning problem, not an infinite renormalization.
This is a logical improvement but not a complete solution.
\end{remark}

%=====================================================================
\section{Open Problems}
%=====================================================================

\begin{enumerate}
\item \textbf{SM particle masses.} The 600-cell adjacency matrix has
  exactly six distinct eigenvalues, mapping suggestively to three SM
  generations with two members each. No derivation of mass ratios from
  the golden-ratio eigenvalues yet exists.

\item \textbf{Fine-structure constant.} $\alpha = 1/137$ has not been
  derived from sea\_strength and the 600-cell geometry.

\item \textbf{Quantum gravity.} The Nexus near Planck-scale horizons
  and the lattice interpretation of Hawking radiation remain open.

\item \textbf{Born rule from first principles.} Theorem~\ref{thm:born}
  borrows from companion~C3. A derivation from the four primitives
  of Definition~\ref{def:primitives} alone would make the series
  fully self-contained.

\item \textbf{Pointer basis test.} The prediction that the pointer basis
  is fixed by the lattice (not the apparatus) has no proposed
  experimental realisation.
\end{enumerate}

%=====================================================================
\section{Consolidated Predictions}
%=====================================================================

\begin{longtable}{@{}p{4.5cm}p{4.5cm}p{3.5cm}@{}}
\caption{Falsifiable predictions of CPP quantum mechanics.}
\label{tab:predictions} \\
\toprule
Prediction & Formula & Status \\
\midrule
\endfirsthead
\toprule
Prediction & Formula & Status \\
\midrule
\endhead
\bottomrule
\endfoot
Bell correlations
& $E = -\cos\theta$, $|S| = 2\sqrt{2}$
& Confirmed \citep{hensen2015,giustina2015,shalm2015} \\[4pt]
Lattice dispersion
& $E = h\nu[1-({\nu}/{\nu_P})^2]$
& Planck-scale experiment \\[4pt]
Decoherence rate
& $\gamma = (\text{sea\_str})^2\EP/\hbar$
& Quantum optics timescales \\[4pt]
Pointer basis (CPP-specific)
& SSV eigenstates, lattice-fixed
& Proposed test needed \\[4pt]
Renormalization
& $\Sigma \sim e^2\EP^2$, finite
& Hierarchy: finite fine-tuning \\[4pt]
Six eigenvalues
& $\pm(1{+}\varphi),\pm\varphi,\pm(\varphi{-}1)$
& Open mass-ratio problem \\
\end{longtable}

%=====================================================================
\section{Conclusion}
%=====================================================================

The five theorems presented here derive, from a single discrete lattice
structure, the Schr\"{o}dinger equation, the Born rule, Bell inequality
violation at the Tsirelson bound, the Lindblad decoherence equation
with a lattice-fixed pointer basis, and bosonic/fermionic commutation
relations. Each theorem identifies the specific physical ingredient:
complex hopping (imaginary unit from phase), non-separable Nexus states
(entanglement), DP~Sea scattering (measurement), and CP exclusion
(fermions). The framework makes one genuinely novel prediction
distinguishable from standard quantum mechanics: the pointer basis is
fixed by the lattice geometry rather than by the measurement apparatus,
and does not depend on the specific detector.

%=====================================================================
\appendix
\section{Graph Laplacian of the 600-Cell}
\label{app:laplacian}
%=====================================================================

The 600-cell has $z=12$ nearest neighbours in $d=3$ dimensions.
For a smooth $\psi(\mathbf{r})$, icosahedral symmetry implies
$\sum_{j\sim i}(\mathbf{r}_j - \mathbf{r}_i) = 0$ and
$\sum_{j\sim i}(\mathbf{r}_j - \mathbf{r}_i)\otimes(\mathbf{r}_j - \mathbf{r}_i)
= (z\Delta s^2/d)\mathbf{I}$, giving:
\begin{equation}
\sum_{j\sim i}(\psi_j - \psi_i)
= \frac{z\Delta s^2}{2d}\nabla^2\psi + O(\Delta s^4)
= 2\Delta s^2\nabla^2\psi + O(\Delta s^4).
\label{eq:laplacian_proof}
\end{equation}
The numerical value $z/(2d) = 12/6 = 2$ is a property of the 600-cell
coordination number. This is why the correct hopping amplitude is
$T = \hbar^2/(4m\Delta s^2)$ — the factor of 4 absorbs the factor of 2
from the Laplacian.

%=====================================================================
\bibliographystyle{unsrtnat}
\begin{thebibliography}{20}

\bibitem{cpp2040a}
T.~L. Abshier and Grok,
``Schr\"{o}dinger equation from phase-carrying DI-bit hopping,''
CPP-2040a v3.1, Hyperphysics Institute (2026).

\bibitem{cpp2040b}
T.~L. Abshier and Grok,
``Superposition and interference from multi-path DI-bit summation,''
CPP-2040b-super v3.1, Hyperphysics Institute (2026).

\bibitem{cpp2040c}
T.~L. Abshier and Grok,
``Entanglement and Bell violation from non-separable Nexus states,''
CPP-2040c v3.1, Hyperphysics Institute (2026).

\bibitem{cpp2040d}
T.~L. Abshier and Grok,
``Measurement problem from DP~Sea decoherence,''
CPP-2040d v3.1, Hyperphysics Institute (2026).

\bibitem{cpp2040e}
T.~L. Abshier and Grok,
``Emergent QFT and finite renormalization from the 600-cell lattice,''
CPP-2040e v3.1, Hyperphysics Institute (2026).

\bibitem{cpp5014}
T.~L. Abshier and Grok,
``Charge neutrality in baryons and emergent quark charges,''
CPP-5014, Hyperphysics Institute (2026).

\bibitem{abshier2026c1}
T.~L. Abshier,
``The Absolute Moment Postulate,''
companion~C1, Hyperphysics Institute (2026).

\bibitem{abshier2026c3}
T.~L. Abshier,
``The Born rule from DI-bit density,''
companion~C3, Hyperphysics Institute (2026).

\bibitem{bell1964}
J.~S. Bell,
``On the Einstein Podolsky Rosen paradox,''
\textit{Physics Physique Fizika} \textbf{1}(3), 195--200 (1964).

\bibitem{clauser1969}
J.~F. Clauser, M.~A. Horne, A.~Shimony, and R.~A. Holt,
``Proposed experiment to test local hidden-variable theories,''
\textit{Phys.\ Rev.\ Lett.} \textbf{23}, 880 (1969).

\bibitem{hensen2015}
B.~Hensen \textit{et al.},
``Loophole-free Bell inequality violation,''
\textit{Nature} \textbf{526}, 682 (2015).

\bibitem{giustina2015}
M.~Giustina \textit{et al.},
``Significant-loophole-free test of Bell's theorem,''
\textit{Phys.\ Rev.\ Lett.} \textbf{115}, 250401 (2015).

\bibitem{shalm2015}
L.~K. Shalm \textit{et al.},
``Strong loophole-free test of local realism,''
\textit{Phys.\ Rev.\ Lett.} \textbf{115}, 250402 (2015).

\bibitem{zurek1981}
W.~H. Zurek,
``Pointer basis of quantum apparatus,''
\textit{Phys.\ Rev.\ D} \textbf{24}, 1516 (1981).

\bibitem{lindblad1976}
G.~Lindblad,
``On the generators of quantum dynamical semigroups,''
\textit{Commun.\ Math.\ Phys.} \textbf{48}, 119 (1976).

\bibitem{brouwer1989}
A.~E. Brouwer, A.~M. Cohen, and A.~Neumaier,
\textit{Distance-Regular Graphs}, Springer (1989).

\bibitem{feynman1965}
R.~P. Feynman and A.~R. Hibbs,
\textit{Quantum Mechanics and Path Integrals}, McGraw-Hill (1965).

\bibitem{breuer2002}
H.-P. Breuer and F.~Petruccione,
\textit{The Theory of Open Quantum Systems}, Oxford (2002).

\end{thebibliography}

\end{document}
